%% kapa: exam=1 round=23 year=2012 period=2 form=A subject=appraisal_law subject_ko=부동산관계법규 seq=24 number=64
%% source: Q-Net 감정평가사 기출문제 (cst00309) artlSeq=5052580 file=raw/1008/5052580/2012년 제23회 감정평가사 제1차시험 기출문제(2교시 A형).pdf
%% conversion: kordoc 4.24.51 -> tools/gen1.py (auto); review=proofread-sonnet; flags=none
\begin{question}{64}
국토해양부장관은 성실의무 등의 위반을 이유로 A 감정평가법인에 대한 업무정지 6개월의 처분에 갈음하여 과징금을 부과하고자 한다. 이 경우 부동산 가격공시 및 감정평가에 관한 법령상 국토해양부장관이 위반행위와 관련한 제반 사정을 종합적으로 고려하여 부과할 수 있는 과징금의 최저액수는?
\begin{choices}
  \item 5천만 원
  \item 1억 원
  \item 1억 2천 5백만 원
  \item 1억 7천 5백만 원
  \item 2억 5천만 원
\end{choices}
\end{question}
\answer{3}
